\section*{2. 第 XVIII 讲的一个补充}

\subsubsection*{命题 2.1.}
设 $X\subset\mathbb P$ 为代数闭域 $k$ 上纯奇维数 $2n+1$ 的光滑射影簇, $\ell\subset\check{\mathbb P}$ 为 $X$ 的超平面截面的泛 Lefschetz 线束. 设 $s$ 为 $\ell$ 的几何点, $H_s$ 为对应的超平面截面, $S\subset\ell$ 为线束的例外值集, $\pi_1(\ell-S,s)$ 为单值群. 假设 $X$ 满足条件 (LV). 则 $\pi_1(\ell-S,s)$ 在消灭上同调
\[
\operatorname{Ev}(H_s,\mathbb Q_\ell)
\subset H^{2n}(H_s,\mathbb Q_\ell)
\]
上的表示绝对不可约.\footnote{原印本说 $X$ 的维数为偶数 $2n$. 上述显示式中的次数 $H^{2n}(H_s)$、反射公式、注 2.2 以及对定理 1.4 的应用均要求维数为奇数 $2n+1$.}

回顾, $\operatorname{Ev}(H_s,\mathbb Q_\ell)$ 是 $H^{2n}(X,\mathbb Q_\ell)$ 在 $H^{2n}(H_s,\mathbb Q_\ell)$ 中的像的正交补, 而此像等于单值不变部分（XVIII 5.6.10）.

\subsubsection*{证明.}
射影直线 $\mathbb P^1$ 是单连通的. 因此 $\pi_1(\ell-S,s)$ 在拓扑意义下由 $S$ 中各点处的惯性子群及其共轭生成. 从而它在 $\operatorname{GL}(H^{2n}(H_s,\mathbb Q_\ell))$ 中的像, 在拓扑意义下由变换
\[
x\longmapsto x-(-1)^n(x\cdot\delta)\delta,
\]
生成, 其中 $\delta$ 共轭于某个消灭圈.\footnote{原印本此处写作 $H^n(H_s)$, 虽然其上下文中的所有出现以及命题的陈述均使用中间次数 $H^{2n}(H_s)$.}
因此, $H^{2n}(X,\mathbb Q_\ell)$ 在 $H^{2n}(H_s,\mathbb Q_\ell)$ 中的像, 是消灭圈及其共轭所生成子空间的正交补; 从而 $\operatorname{Ev}(H_s,\mathbb Q_\ell)$ 由这些圈及其共轭生成.

对一个\underline{泛} Lefschetz 线束, 第 XVIII 讲第 6 节的论证说明, $S$ 的不同点所对应的消灭圈彼此共轭.

设 $W\subset\operatorname{Ev}(H_s,\mathbb Q_\ell)$ 为非零不变子空间. 存在某个消灭圈的共轭 $\delta$ 以及元素 $v\in W$, 使 $(v,\delta)\ne0$; 这里用到了条件 (LV). 因为
\[
v-(-1)^n(v\cdot\delta)\delta\in W,
\]
所以 $\delta\in W$.\footnote{原印本在第二次出现的这一式中使用了未定义的指数 $m$; 紧前的反射公式及 Picard--Lefschetz 公式给出的指数为 $n$.} 子空间 $W$ 包含一个消灭圈及其所有共轭, 因而等于 $\operatorname{Ev}(H_s,\mathbb Q_\ell)$. 此论证在扩张系数域 $\mathbb Q_\ell$ 后仍然成立, 从而证明了绝对不可约性.

\subsubsection*{注 2.2.}
若去掉条件 (LV), 上述论证仍说明, $H^{2n}(H_s,\mathbb Q_\ell)$ 的每个不变子空间, 或者包含 $\operatorname{Ev}(H_s,\mathbb Q_\ell)$, 或者包含于 $H^{2n}(X,\mathbb Q_\ell)$ 的像中. 对偶维数的 $X$（以及任意 Lefschetz 线束）, 同样的断言成立.
